% !TeX root = main.tex
\IfFileExists{../apunte_mecanica_clasica.tex}{%
  \documentclass[../apunte_mecanica_clasica.tex]{subfiles}%
}{%
  \documentclass[apunte_mecanica_clasica.tex]{subfiles}%
}

\begin{document}

\chapter{Métrica}

\section{Longitud de trayectoria y elemento de línea}

Una partícula al moverse en $\mathbb{R}^3$ describe una curva, que 
llamamos \textbf{Trayectoria}, y cada punto está descrito por un
vector posición. Puesto que la trayectoria es continua, 
esta se concibe como un conjunto continuo de posiciones.

Una de las magnitudes físicas de interés es medir la 
\textbf{longitud de la trayectoria}, esto es, la 
\textbf{distancia recorrida} por la partícula.

Recurramos al cálculo diferencial: consideremos un trozo de la trayectoria 
entre dos puntos infinitesimalmente vecinos y aproximémoslo por una línea 
recta, tal que 
\begin{equation*}
    ds = \lVert d\vec{x} \rVert \,,
\end{equation*}
y, por simplicidad, usaremos las coordenadas cartesianas.

\begin{figure}[ht]
    \centering
    \includegraphics[width=0.6\linewidth]{figures/04_metrica/met_fig1.png}
    \caption{A la izquierda: el elemento infinitesimal de longitud $ds$, 
    con los vectores posición $\vec{x}$ y $\vec{x} + d\vec{x}$, unidos 
    por el vector $d\vec{x}$. A la derecha: representación del vector 
    $d\vec{x}$ en el espacio, cuya norma es igual a $ds$, junto con 
    la proyección de sus componentes sobre los ejes cartesianos $x$, $y$ 
    y $z$.}
\end{figure}

La longitud del segmento infinitesimal de línea, $ds$, puede ser 
encontrada usando el teorema de Pitágoras: 
\begin{equation*}
    ds^2 = dT^2 + dz^2 = dx^2 + dy^2 + dz^2 \,,
\end{equation*}
es decir,
\begin{equation*}
    ds^2 = dx^2 + dy^2 + dz^2 \,.
\end{equation*}

Esto corresponde a la distancia infinitesimal en coordenadas cartesianas. 
Por lo tanto, la distancia total de la trayectoria es
\begin{equation*}
    \Delta S = \int_{\vec{x}_i}^{\vec{x}_f} 
    \sqrt{dx^2 + dy^2 + dz^2} \,.
\end{equation*}

\section{Elemento de línea en coordenadas esféricas}

El mismo elemento de línea infinitesimal también puede ser 
calculado en otros sistemas coordenados, por ejemplo, en 
coordenadas esféricas. 

Recordemos que las líneas coordenadas esféricas son:
\begin{itemize}
    \item \textbf{Líneas $r$:} rectas que pasan por el origen; $\theta$ y $\phi$ constantes, $r$ variable.
    \item \textbf{Líneas $\phi$:} circunferencias horizontales de radio $r\sin\theta$ (a altura $z = r\cos\theta$); $r$ y $\theta$ constantes, $\phi$ variable.
    \item \textbf{Líneas $\theta$:} semicírculos en planos meridianos (que contienen al eje $z$) de radio $r$; $r$ y $\phi$ constantes, $\theta$ variable.
\end{itemize}

\begin{figure}[h]
    \centering
    \includegraphics[width=0.55\linewidth]{figures/04_metrica/met_fig2.png}
    \caption{Elemento infinitesimal en coordenadas esféricas. 
    Se muestran los incrementos diferenciales: 
    radial $dr$, polar $r\,d\theta$ y azimutal $r\sin\theta\,d\phi$, 
    que en conjunto definen el diferencial de línea $ds$ en este sistema.}
\end{figure}

Puesto que las líneas coordenadas son ortogonales, se tiene que
\begin{align*}
    ds^2 &= (ds_1)^2 + (ds_2)^2 + (ds_3)^2 \\[2pt]
    &= (dr)^2 + (r \sin\theta \, d\phi)^2 + (r \, d\theta)^2 \,.
\end{align*}

Es decir,
\begin{equation*}
    ds^2 = dr^2 + r^2 \sin^2\theta \, d\phi^2 + r^2 d\theta^2 \,,
\end{equation*}
lo que corresponde a la distancia entre dos puntos infinitesimalmente 
cercanos en coordenadas esféricas, y debe ser igual a la encontrada 
en coordenadas cartesianas:
\begin{align*}
    ds^2 &= dx^2 + dy^2 + dz^2 \\[2pt]
         &= dr^2 + r^2 \sin^2\theta \, d\phi^2 + r^2 d\theta^2 \,.
\end{align*}

\section{Elemento de línea en coordenadas generalizadas}

La cantidad $ds^2$ se llama \textbf{Elemento de Línea}. 
Ahora escribamos el elemento de línea en coordenadas generalizadas. 

Usando
\begin{equation*}
    d\vec{x} = dx \, \hat{\imath} + dy \, \hat{\jmath} + dz \, \hat{k} \,,
\end{equation*}
entonces
\begin{align*}
    ds^2 &= dx^2 + dy^2 + dz^2 \\[2pt]
          &= dx \, \hat{\imath} \cdot dx \, \hat{\imath} 
           + dy \, \hat{\jmath} \cdot dy \, \hat{\jmath} 
           + dz \, \hat{k} \cdot dz \, \hat{k} \\[2pt]
          &= (dx \, \hat{\imath} + dy \, \hat{\jmath} + dz \, \hat{k})
             \cdot (dx \, \hat{\imath} + dy \, \hat{\jmath} + dz \, \hat{k}) \,.
\end{align*}

Por lo tanto,
\begin{equation}\label{4_metrica_ds^2}
    ds^2 = d\vec{x} \cdot d\vec{x} \,.
\end{equation}

En un sistema de coordenadas $Ox^1x^2x^3$, la cantidad $d\vec{x}$ se escribe como
\begin{align*}
    d\vec{x}(x^1, x^2, x^3) &= 
    \frac{\partial \vec{x}}{\partial x^1} \, dx^1 
    + \frac{\partial \vec{x}}{\partial x^2} \, dx^2 
    + \frac{\partial \vec{x}}{\partial x^3} \, dx^3 \\[4pt]
    &= \vec{e}_1 \, dx^1 + \vec{e}_2 \, dx^2 + \vec{e}_3 \, dx^3 \,.
\end{align*}

De forma compacta:
\begin{equation*}
    d\vec{x} = \sum_{i=1}^3 \vec{e}_i \, dx^i \,.
\end{equation*}

Reemplazando en \eqref{4_metrica_ds^2}, se obtiene:
\begin{align}
    ds^2 &= \left( \sum_{i=1}^3 \vec{e}_i \, dx^i \right) 
            \cdot 
            \left( \sum_{j=1}^3 \vec{e}_j \, dx^j \right) \nonumber \\[2pt]
          &= \sum_{i=1}^3 \sum_{j=1}^3 
             \big( \vec{e}_i \cdot \vec{e}_j \big) \, dx^i dx^j \,.
             \label{4_metrica_elemento_linea_general}
\end{align}

\section{Coeficientes métricos}

\begin{apuntebloque}{Definición}
Se definen los \textbf{coeficientes métricos} (o simplemente la 
\textbf{métrica}), asociados al sistema coordenado $Ox^1x^2x^3$, como:
\begin{equation}\label{4_metrica_g}
    g_{ij} = \vec{e}_i \cdot \vec{e}_j \,.
\end{equation}
\end{apuntebloque}

Por lo tanto, existen $3 \times 3 = 9$ coeficientes métricos. 
Así, el elemento de línea \eqref{4_metrica_elemento_linea_general} toma la forma
\begin{equation}\label{4_metrica_elemento_linea_sistema_coordenado_generalizado}
    ds^2 = \sum_{i=1}^3 \sum_{j=1}^3 g_{ij} \, dx^i dx^j \,.
\end{equation}

Esta expresión corresponde al elemento de línea escrito en el 
sistema coordenado $Ox^1x^2x^3$ generalizado.

\section{Ejemplos fundamentales}

\subsection{Coordenadas cartesianas}

\begin{apuntebloque}{Ejemplo: Coordenadas Cartesianas}
En \textbf{coordenadas cartesianas}:
\begin{equation*}
    x^1 = x,\quad x^2 = y,\quad x^3 = z, \qquad
    \vec{e}_1=\hat{\imath},\ \vec{e}_2=\hat{\jmath},\ \vec{e}_3=\hat{k}.
\end{equation*}

Coeficientes métricos:
\begin{align*}
    g_{11}&=\vec{e}_1\!\cdot\!\vec{e}_1=1,&
    g_{12}&=\vec{e}_1\!\cdot\!\vec{e}_2=0,&
    g_{13}&=\vec{e}_1\!\cdot\!\vec{e}_3=0,\\
    g_{21}&=\vec{e}_2\!\cdot\!\vec{e}_1=0,&
    g_{22}&=\vec{e}_2\!\cdot\!\vec{e}_2=1,&
    g_{23}&=\vec{e}_2\!\cdot\!\vec{e}_3=0,\\
    g_{31}&=\vec{e}_3\!\cdot\!\vec{e}_1=0,&
    g_{32}&=\vec{e}_3\!\cdot\!\vec{e}_2=0,&
    g_{33}&=\vec{e}_3\!\cdot\!\vec{e}_3=1.
\end{align*}

Elemento de línea:
\begin{align*}
    ds^2
      &= \sum_{i=1}^3\sum_{j=1}^3 g_{ij}\,dx^i dx^j \\[2pt]
      &= \sum_{j=1}^3 g_{1j}\,dx^1 dx^j
       + \sum_{j=1}^3 g_{2j}\,dx^2 dx^j
       + \sum_{j=1}^3 g_{3j}\,dx^3 dx^j \\[2pt]
      &= g_{11}(dx^1)^2 + g_{12}\,dx^1 dx^2 + g_{13}\,dx^1 dx^3 \\
      &\quad + g_{21}\,dx^2 dx^1 + g_{22}(dx^2)^2 + g_{23}\,dx^2 dx^3 \\
      &\quad + g_{31}\,dx^3 dx^1 + g_{32}\,dx^3 dx^2 + g_{33}(dx^3)^2 \\[2pt]
      &= 1\cdot(dx^1)^2 + 0\cdot dx^1 dx^2 + 0\cdot dx^1 dx^3 \\
      &\quad + 0\cdot dx^2 dx^1 + 1\cdot(dx^2)^2 + 0\cdot dx^2 dx^3 \\
      &\quad + 0\cdot dx^3 dx^1 + 0\cdot dx^3 dx^2 + 1\cdot(dx^3)^2 \\[4pt]
      &= (dx^1)^2 + (dx^2)^2 + (dx^3)^2 \\[2pt]
      &= dx^2 + dy^2 + dz^2 \,.
\end{align*}
\begin{equation*}
    \therefore \quad ds^2=dx^2 + dy^2 + dz^2 \,.
\end{equation*}
En consecuencia, la métrica cartesiana es la \textbf{matriz identidad} de orden $3$:
\begin{equation*}
    [g_{ij}] = I_3 \qquad (\text{equivalentemente, } g_{ij}=\delta_{ij}).
\end{equation*}
\end{apuntebloque}

\clearpage

\subsection{Coordenadas esféricas}

\begin{apuntebloque}{Ejemplo: Coordenadas Esféricas}
En \textbf{coordenadas esféricas} tomamos
\begin{equation*}
    x^1 = r, \qquad x^2 = \phi, \qquad x^3 = \theta \,.
\end{equation*}

Las bases coordenadas (no unitarias en general) son
\begin{align*}
    \vec{e}_r(\phi,\theta) &= \cos\phi\,\sin\theta\,\hat{\imath}
                           + \sin\phi\,\sin\theta\,\hat{\jmath}
                           + \cos\theta\,\hat{k},\\
    \vec{e}_\phi(r,\phi,\theta) &= r\sin\theta\,
        \big(-\sin\phi\,\hat{\imath}+\cos\phi\,\hat{\jmath}\big),\\
    \vec{e}_\theta(r,\phi,\theta) &= r\,
        \big(\cos\phi\cos\theta\,\hat{\imath}
            +\sin\phi\cos\theta\,\hat{\jmath}
            -\sin\theta\,\hat{k}\big).
\end{align*}

Por lo tanto, los coeficientes métricos son
\begin{align*}
    g_{11} &= \vec{e}_r\!\cdot\!\vec{e}_r = 1, &
    g_{12} &= \vec{e}_r\!\cdot\!\vec{e}_\phi = 0, &
    g_{13} &= \vec{e}_r\!\cdot\!\vec{e}_\theta = 0,\\
    g_{21} &= \vec{e}_\phi\!\cdot\!\vec{e}_r = 0, &
    g_{22} &= \vec{e}_\phi\!\cdot\!\vec{e}_\phi = r^2\sin^2\theta, &
    g_{23} &= \vec{e}_\phi\!\cdot\!\vec{e}_\theta = 0,\\
    g_{31} &= \vec{e}_\theta\!\cdot\!\vec{e}_r = 0, &
    g_{32} &= \vec{e}_\theta\!\cdot\!\vec{e}_\phi = 0, &
    g_{33} &= \vec{e}_\theta\!\cdot\!\vec{e}_\theta = r^2.
\end{align*}

Elemento de línea (expansión):
\begin{align*}
    ds^2
      &= \sum_{i=1}^3\sum_{j=1}^3 g_{ij}\,dx^i dx^j \\[2pt]
      &= \sum_{j=1}^3 g_{1j}\,dx^1 dx^j
       + \sum_{j=1}^3 g_{2j}\,dx^2 dx^j
       + \sum_{j=1}^3 g_{3j}\,dx^3 dx^j \\[2pt]
      &= g_{11}(dx^1)^2 + g_{12}\,dx^1 dx^2 + g_{13}\,dx^1 dx^3 \\
      &\quad + g_{21}\,dx^2 dx^1 + g_{22}(dx^2)^2 + g_{23}\,dx^2 dx^3 \\
      &\quad + g_{31}\,dx^3 dx^1 + g_{32}\,dx^3 dx^2 + g_{33}(dx^3)^2 \\[4pt]
      &= 1\cdot(dr)^2 + 0\cdot\,dr\,d\phi + 0\cdot\,dr\,d\theta \\
      &\quad + 0\cdot\,d\phi\,dr + (r^2\sin^2\theta)\,(d\phi)^2 + 0\cdot\,d\phi\,d\theta \\
      &\quad + 0\cdot\,d\theta\,dr + 0\cdot\,d\theta\,d\phi + r^2\,(d\theta)^2 \\[4pt]
      &= dr^2 + r^2\sin^2\theta\,d\phi^2 + r^2\,d\theta^2 \,.
\end{align*}
\begin{equation*}
    \therefore \quad ds^2=dr^2 + r^2\sin^2\theta\,d\phi^2 + r^2\,d\theta^2 \,.
\end{equation*}

\end{apuntebloque}

\newpage

\section{Simetría de la métrica y forma matricial}

Notar que la métrica $g_{ij}$ posee la \textbf{propiedad de simetría}:
\begin{align*}
    &g_{ij} \;=\; \vec{e}_i \cdot \vec{e}_j \;=\; \vec{e}_j \cdot \vec{e}_i \;=\; g_{ji} \\[4pt]
    &\therefore \quad g_{ij} = g_{ji} \,.
\end{align*}

Esto es, \textbf{la métrica es invariante bajo el intercambio de índices}. Por lo tanto,
\begin{equation*}
    g_{12} = g_{21}, \qquad
    g_{13} = g_{31}, \qquad
    g_{23} = g_{32}.
\end{equation*}

En consecuencia, el número de coeficientes métricos independientes se reduce de $9$ a $6$.

Para entender esta simetría, interpretemos la métrica como una 
\textbf{matriz de orden $3 \times 3$}. En la expresión $g_{ij}$:
\begin{itemize}
    \item El índice $i$ rotula las filas.
    \item El índice $j$ rotula las columnas.
\end{itemize}

Entonces,
\begin{equation*}
    [g_{ij}] =
    \begin{pmatrix}
        g_{11} & g_{12} & g_{13} \\
        g_{21} & g_{22} & g_{23} \\
        g_{31} & g_{32} & g_{33}
    \end{pmatrix}
    =
    \begin{pmatrix}
        \vec{e}_1 \cdot \vec{e}_1 & \vec{e}_1 \cdot \vec{e}_2 & \vec{e}_1 \cdot \vec{e}_3 \\
        \vec{e}_2 \cdot \vec{e}_1 & \vec{e}_2 \cdot \vec{e}_2 & \vec{e}_2 \cdot \vec{e}_3 \\
        \vec{e}_3 \cdot \vec{e}_1 & \vec{e}_3 \cdot \vec{e}_2 & \vec{e}_3 \cdot \vec{e}_3
    \end{pmatrix}
    = g \,.
\end{equation*}

Por lo tanto, $g_{ij} = g_{ji}$ equivale al \textbf{intercambio de filas por 
columnas}, lo que en lenguaje matricial se escribe como
\begin{equation*}
    g^t = g \,,
\end{equation*}
donde $t$ significa \textbf{transpuesta}. Es decir, la matriz $g$ es 
\textbf{simétrica}.

Entonces, el triángulo superior es igual al triángulo inferior en la 
matriz métrica. De los seis elementos de los dos triángulos de la matriz, 
solo tres son independientes, y si a estos se les suman los tres elementos 
de la diagonal, se concluye que la matriz métrica posee únicamente 
\textbf{seis elementos independientes}.

\section{Elemento de línea en forma matricial}

El elemento de línea, aunque sea un escalar, se puede escribir en términos 
matriciales. Escribamos los diferenciales $dx^i$, $i\in\{1,2,3\}$, como una 
\textbf{matriz columna de orden $3 \times 1$}:
\begin{equation*}
    dx^i \;\;\longrightarrow\;\;
    dX =
    \begin{pmatrix}
        dx^1 \\
        dx^2 \\
        dx^3
    \end{pmatrix}.
\end{equation*}

Así, el elemento de línea se escribe como
\begin{equation*}
    ds^2 = dX^t \, g \, dX \,,
\end{equation*}
donde $dX^t$ es la transpuesta de $dX$.

\begin{apuntebloque}{Ejemplo: Métrica en diferentes coordenadas}

\textbf{a) Coordenadas Cartesianas}

La métrica se obtiene como
\begin{equation*}
    g_{ij} =
    \begin{pmatrix}
        \hat{\imath}\!\cdot\!\hat{\imath} & \hat{\imath}\!\cdot\!\hat{\jmath} & \hat{\imath}\!\cdot\!\hat{k} \\
        \hat{\jmath}\!\cdot\!\hat{\imath} & \hat{\jmath}\!\cdot\!\hat{\jmath} & \hat{\jmath}\!\cdot\!\hat{k} \\
        \hat{k}\!\cdot\!\hat{\imath} & \hat{k}\!\cdot\!\hat{\jmath} & \hat{k}\!\cdot\!\hat{k}
    \end{pmatrix}
    =
    \begin{pmatrix}
        1 & 0 & 0 \\
        0 & 1 & 0 \\
        0 & 0 & 1
    \end{pmatrix}.
\end{equation*}

Esta es la \textbf{métrica asociada a las coordenadas cartesianas}, y se denomina
\textbf{Métrica Euclidiana}.  

El elemento de línea se escribe en forma matricial:
\begin{align*}
    ds^2 &= 
    \begin{pmatrix}
        dx & dy & dz
    \end{pmatrix}
    \begin{pmatrix}
        1 & 0 & 0 \\
        0 & 1 & 0 \\
        0 & 0 & 1
    \end{pmatrix}
    \begin{pmatrix}
        dx \\ dy \\ dz
    \end{pmatrix} \\[4pt]
    &= 
    \begin{pmatrix}
        dx & dy & dz
    \end{pmatrix}
    \begin{pmatrix}
        dx \\ dy \\ dz
    \end{pmatrix} \\[4pt]
    &= dx^2 + dy^2 + dz^2 \,.
\end{align*}

\textbf{b) Coordenadas Esféricas}

La métrica se obtiene como
\begin{equation*}
    g_{ij} =
    \begin{pmatrix}
        \vec{e}_r\!\cdot\!\vec{e}_r & \vec{e}_r\!\cdot\!\vec{e}_\phi & \vec{e}_r\!\cdot\!\vec{e}_\theta \\
        \vec{e}_\phi\!\cdot\!\vec{e}_r & \vec{e}_\phi\!\cdot\!\vec{e}_\phi & \vec{e}_\phi\!\cdot\!\vec{e}_\theta \\
        \vec{e}_\theta\!\cdot\!\vec{e}_r & \vec{e}_\theta\!\cdot\!\vec{e}_\phi & \vec{e}_\theta\!\cdot\!\vec{e}_\theta
    \end{pmatrix}
    =
    \begin{pmatrix}
        1 & 0 & 0 \\
        0 & r^2\sin^2\theta & 0 \\
        0 & 0 & r^2
    \end{pmatrix}.
\end{equation*}

Esta es la \textbf{métrica asociada a las coordenadas esféricas}.  

El elemento de línea se escribe en forma matricial:
\begin{align*}
    ds^2 &=
    \begin{pmatrix}
        dr & d\phi & d\theta
    \end{pmatrix}
    \begin{pmatrix}
        1 & 0 & 0 \\
        0 & r^2\sin^2\theta & 0 \\
        0 & 0 & r^2
    \end{pmatrix}
    \begin{pmatrix}
        dr \\ d\phi \\ d\theta
    \end{pmatrix} \\[4pt]
    &= 
    \begin{pmatrix}
        dr & d\phi & d\theta
    \end{pmatrix}
    \begin{pmatrix}
        dr \\ r^2\sin^2\theta\,d\phi \\ r^2\,d\theta
    \end{pmatrix} \\[4pt]
    &= dr^2 + r^2\sin^2\theta\,d\phi^2 + r^2\,d\theta^2 \,.
\end{align*}

\bigskip

\end{apuntebloque}

\section{Factores de escala y vectores unitarios}

\begin{apuntebloque}{}

\textbf{Normas de los vectores coordenados}

Se tiene que
\begin{align*}
    \|\vec{e}_1\| &= \sqrt{\vec{e}_1 \cdot \vec{e}_1} = \sqrt{g_{11}},\\
    \|\vec{e}_2\| &= \sqrt{\vec{e}_2 \cdot \vec{e}_2} = \sqrt{g_{22}},\\
    \|\vec{e}_3\| &= \sqrt{\vec{e}_3 \cdot \vec{e}_3} = \sqrt{g_{33}}.
\end{align*}

De manera compacta:
\begin{equation*}
    \|\vec{e}_i\| = \sqrt{g_{ii}}, 
    \qquad i \in \{1,2,3\}.
\end{equation*}

Por lo tanto, los vectores unitarios se escriben como
\begin{equation*}
    \hat{e}_i = \frac{\vec{e}_i}{\sqrt{g_{ii}}},
    \qquad i \in \{1,2,3\}.
\end{equation*}

Notar que
\begin{equation*}
    h_i = \|\vec{e}_i\| = \sqrt{g_{ii}}, 
    \qquad i \in \{1,2,3\},
\end{equation*}
son los llamados \textbf{factores de escala} cuando las 
coordenadas son ortogonales.

\end{apuntebloque}

\section{Ortogonalidad de líneas coordenadas}

La definición de la métrica involucra a los vectores base 
$\{\vec{e}_i\}_{i=1}^3$, y no a los vectores unitarios. 
Por lo tanto, la métrica contiene la información propia 
del sistema coordenado $Ox^1x^2x^3$, en particular 
dónde se anula, diverge, etc.

\begin{apuntebloque}{Proposición}
Las líneas coordenadas generadas por las coordenadas $x^i$ y $x^j$, 
con $i \neq j$, son ortogonales si y solo si
\begin{equation*}
    g_{ij} = 0, \qquad i \neq j \,.
\end{equation*}
\end{apuntebloque}

\newpage

\section{Delta de Kronecker y métrica euclideana}

La métrica Euclideana puede ser escrita en una notación más 
compacta definiendo:

\begin{apuntebloque}{Definición}
Se define la \textbf{delta de Kronecker} como
\begin{equation*}
\delta_{ij} =
\begin{cases}
    1, & \text{si } i = j, \\[6pt]
    0, & \text{si } i \neq j.
\end{cases}
\end{equation*}
\end{apuntebloque}

Escribiendo la delta de Kronecker en forma explícita:
\begin{align*}
    \delta_{11} &= 1, & \delta_{12} &= 0, & \delta_{13} &= 0, \\
    \delta_{21} &= 0, & \delta_{22} &= 1, & \delta_{23} &= 0, \\
    \delta_{31} &= 0, & \delta_{32} &= 0, & \delta_{33} &= 1.
\end{align*}

Escribiéndola en forma matricial:
\begin{equation*}
    [\delta_{ij}] =
    \begin{pmatrix}
        \delta_{11} & \delta_{12} & \delta_{13} \\
        \delta_{21} & \delta_{22} & \delta_{23} \\
        \delta_{31} & \delta_{32} & \delta_{33}
    \end{pmatrix}
    =
    \begin{pmatrix}
        1 & 0 & 0 \\
        0 & 1 & 0 \\
        0 & 0 & 1
    \end{pmatrix}.
\end{equation*}

Por lo tanto, la métrica Euclideana toma la forma compacta:
\begin{equation*}
    g_{ij} = \delta_{ij}.
\end{equation*}

Se debe aclarar que $\delta_{ij}$ no es la matriz identidad, 
sino que es la \textbf{métrica del espacio plano} $\mathbb{R}^3$ 
descrito en coordenadas cartesianas.  
De manera análoga, la métrica en coordenadas esféricas también 
describe el espacio plano $\mathbb{R}^3$, pero en coordenadas 
esféricas.  
Como veremos más adelante, ambas métricas están relacionadas 
por una \textbf{transformación de coordenadas}.

\clearpage

Sean $\vec{A}$ y $\vec{B}$ dos vectores descritos en el sistema de 
coordenadas $Ox^1x^2x^3$, con base coordenada 
$\{\vec{e}_i\}_{i=1}^3$. El producto escalar entre ellos toma la forma:
\begin{align*}
    \vec{A} \cdot \vec{B} 
        &= \left( \sum_{i=1}^3 A^i \vec{e}_i \right) 
           \cdot 
           \left( \sum_{j=1}^3 B^j \vec{e}_j \right) \\[4pt]
        &= \sum_{i=1}^3 \sum_{j=1}^3 
           \big( \vec{e}_i \cdot \vec{e}_j \big) A^i B^j \\[4pt]
        &= \sum_{i=1}^3 \sum_{j=1}^3 g_{ij} \, A^i B^j \,.
\end{align*}

\begin{apuntebloque}{Definición}
El \textbf{producto escalar} entre los vectores $\vec{A}$ y $\vec{B}$, 
en el sistema coordenado $Ox^1x^2x^3$, es
\begin{equation*}
    \vec{A} \cdot \vec{B} 
    = \sum_{i=1}^3 \sum_{j=1}^3 g_{ij}\,A^i B^j \,.
\end{equation*}
\end{apuntebloque}

\begin{apuntebloque}{Definición}
La \textbf{magnitud} de un vector $\vec{A}$, en el sistema coordenado 
$Ox^1x^2x^3$, es
\begin{equation*}
    \|\vec{A}\| = \sqrt{\vec{A} \cdot \vec{A}}
    = \sqrt{\sum_{i=1}^3 \sum_{j=1}^3 g_{ij}\,A^i A^j} \,.
\end{equation*}
\end{apuntebloque}

\begin{apuntebloque}{Ejemplo: Producto escalar y magnitud en distintos sistemas coordenados}

\textbf{a) Coordenadas Cartesianas}

\begin{align*}
    \vec{A} \cdot \vec{B} 
        &= \sum_{i=1}^3 \sum_{j=1}^3 g_{ij} A^i B^j \\[4pt] 
        &= \sum_{i=1}^3 \sum_{j=1}^3 \delta_{ij} A^i B^j \\[4pt]
        &= \sum_{j=1}^3 \delta_{1j} A^1 B^j 
         + \sum_{j=1}^3 \delta_{2j} A^2 B^j 
         + \sum_{j=1}^3 \delta_{3j} A^3 B^j \\[4pt]
        &= \delta_{11}A^1B^1 + \delta_{12}A^1B^2 + \delta_{13}A^1B^3 \\
        &\quad + \delta_{21}A^2B^1 + \delta_{22}A^2B^2 + \delta_{23}A^2B^3 \\
        &\quad + \delta_{31}A^3B^1 + \delta_{32}A^3B^2 + \delta_{33}A^3B^3 \\[4pt]
        &= A^1B^1 + A^2B^2 + A^3B^3 \,.
\end{align*}

Por lo tanto,
\begin{equation*}
    \vec{A} \cdot \vec{B} = A^1B^1 + A^2B^2 + A^3B^3 \,.
\end{equation*}

La magnitud del vector $\vec{A}$ es
\begin{equation*}
    \|\vec{A}\| = \sqrt{(A^1)^2 + (A^2)^2 + (A^3)^2} \,.
\end{equation*}

\bigskip

\end{apuntebloque}

\begin{apuntebloque}{}

\textbf{b) Coordenadas Esféricas}

\begin{align*}
    \vec{A} \cdot \vec{B} 
        &= \sum_{i=1}^3 \sum_{j=1}^3 g_{ij} A^i B^j \\[4pt]
        &= \sum_{j=1}^3 g_{1j} A^r B^j 
         + \sum_{j=1}^3 g_{2j} A^\phi B^j
         + \sum_{j=1}^3 g_{3j} A^\theta B^j \\[4pt]
        &= g_{11}A^r B^r + g_{12}A^r B^\phi + g_{13}A^r B^\theta \\
        &\quad + g_{21}A^\phi B^r + g_{22}A^\phi B^\phi + g_{23}A^\phi B^\theta \\
        &\quad + g_{31}A^\theta B^r + g_{32}A^\theta B^\phi + g_{33}A^\theta B^\theta \\[4pt]
        &= A^r B^r + r^2\sin^2\theta \, A^\phi B^\phi + r^2 A^\theta B^\theta \,.
\end{align*}

Por lo tanto,
\begin{equation*}
    \vec{A} \cdot \vec{B} = A^r B^r + r^2\sin^2\theta \, A^\phi B^\phi + r^2 A^\theta B^\theta \,.
\end{equation*}

La magnitud del vector $\vec{A}$ es
\begin{equation*}
    \|\vec{A}\| = \sqrt{(A^r)^2 + r^2\sin^2\theta \,(A^\phi)^2 + r^2 (A^\theta)^2} \,.
\end{equation*}
\end{apuntebloque}

Si la métrica $g_{ij}$ es vista como una matriz de $3 \times 3$, 
entonces puede tener inversa. Pero, antes de definir la \textbf{métrica inversa}, primero reescribamos 
las operaciones matriciales usando la \textbf{notación de índices}, 
es decir, la \textbf{notación tensorial}.

Sea $A$ una matriz de orden $3 \times 3$ y escribámosla usando 
\textbf{notación de índices}:

\[
A \;\; \longrightarrow \;\; A^i{}_j,
\]

donde:
\begin{itemize}
    \item El índice $i$ rotula las filas.
    \item El índice $j$ rotula las columnas.
\end{itemize}

En forma explícita:
\begin{equation*}
A^i{}_j =
\begin{pmatrix}
    A^1{}_1 & A^1{}_2 & A^1{}_3 \\
    A^2{}_1 & A^2{}_2 & A^2{}_3 \\
    A^3{}_1 & A^3{}_2 & A^3{}_3
\end{pmatrix}.
\end{equation*}

Sean las matrices $A$ y $B$ y sean $\alpha, \beta \in \mathbb{R}$, 
entonces su combinación lineal es
\begin{align*}
    C &= \alpha A + \beta B \\[4pt]
    &=
\alpha
\begin{pmatrix}
    A^1{}_1 & A^1{}_2 & A^1{}_3 \\
    A^2{}_1 & A^2{}_2 & A^2{}_3 \\
    A^3{}_1 & A^3{}_2 & A^3{}_3
\end{pmatrix}
+
\beta
\begin{pmatrix}
    B^1{}_1 & B^1{}_2 & B^1{}_3 \\
    B^2{}_1 & B^2{}_2 & B^2{}_3 \\
    B^3{}_1 & B^3{}_2 & B^3{}_3
\end{pmatrix}\\[4pt]
&=
\begin{pmatrix}
    \alpha A^1{}_1 & \alpha A^1{}_2 & \alpha A^1{}_3 \\
    \alpha A^2{}_1 & \alpha A^2{}_2 & \alpha A^2{}_3 \\
    \alpha A^3{}_1 & \alpha A^3{}_2 & \alpha A^3{}_3
\end{pmatrix}
+
\begin{pmatrix}
    \beta B^1{}_1 & \beta B^1{}_2 & \beta B^1{}_3 \\
    \beta B^2{}_1 & \beta B^2{}_2 & \beta B^2{}_3 \\
    \beta B^3{}_1 & \beta B^3{}_2 & \beta B^3{}_3
\end{pmatrix}\\[4pt]
&=
\begin{pmatrix}
    \alpha A^1{}_1 + \beta B^1{}_1 & \alpha A^1{}_2 + \beta B^1{}_2 & \alpha A^1{}_3 + \beta B^1{}_3 \\
    \alpha A^2{}_1 + \beta B^2{}_1 & \alpha A^2{}_2 + \beta B^2{}_2 & \alpha A^2{}_3 + \beta B^2{}_3 \\
    \alpha A^3{}_1 + \beta B^3{}_1 & \alpha A^3{}_2 + \beta B^3{}_2 & \alpha A^3{}_3 + \beta B^3{}_3
\end{pmatrix}\\[4pt]
&=
\begin{pmatrix}
    C^1{}_1 & C^1{}_2 & C^1{}_3 \\
    C^2{}_1 & C^2{}_2 & C^2{}_3 \\
    C^3{}_1 & C^3{}_2 & C^3{}_3
\end{pmatrix}.
\end{align*}

En notación de índices:
\begin{equation*}
    C^i{}_j = \alpha A^i{}_j + \beta B^i{}_j \,.
\end{equation*}

Ahora consideremos el producto matricial $C = AB$ en notación tensorial:

\begin{align*}
C &= AB \\[4pt] 
&=
\begin{pmatrix}
    A^1{}_1 & A^1{}_2 & A^1{}_3 \\
    A^2{}_1 & A^2{}_2 & A^2{}_3 \\
    A^3{}_1 & A^3{}_2 & A^3{}_3
\end{pmatrix}
\begin{pmatrix}
    B^1{}_1 & B^1{}_2 & B^1{}_3 \\
    B^2{}_1 & B^2{}_2 & B^2{}_3 \\
    B^3{}_1 & B^3{}_2 & B^3{}_3
\end{pmatrix} \\[4pt]
&=
\begin{pmatrix}
    \sum_{k=1}^3 A^1{}_k B^k{}_1 & \sum_{k=1}^3 A^1{}_k B^k{}_2 & \sum_{k=1}^3 A^1{}_k B^k{}_3 \\
    \sum_{k=1}^3 A^2{}_k B^k{}_1 & \sum_{k=1}^3 A^2{}_k B^k{}_2 & \sum_{k=1}^3 A^2{}_k B^k{}_3 \\
    \sum_{k=1}^3 A^3{}_k B^k{}_1 & \sum_{k=1}^3 A^3{}_k B^k{}_2 & \sum_{k=1}^3 A^3{}_k B^k{}_3
\end{pmatrix} \\[4pt]
&=
\begin{pmatrix}
    C^1{}_1 & C^1{}_2 & C^1{}_3 \\
    C^2{}_1 & C^2{}_2 & C^2{}_3 \\
    C^3{}_1 & C^3{}_2 & C^3{}_3
\end{pmatrix}.
\end{align*}

En notación de índices:
\begin{equation*}
    C^i{}_j = \sum_{k=1}^3 A^i{}_k B^k{}_j \,.
\end{equation*}

El producto matricial se realiza con un índice arriba y otro abajo 
(índice $k$), esto es porque el producto matricial no debe depender 
del sistema coordenado.  

Un producto matricial de especial interés es el de una matriz y su 
inversa:
\begin{equation*}
    A A^{-1} = A^{-1} A = I \,,
\end{equation*}
donde $I$ es la \textbf{matriz identidad}.  

Por lo tanto, necesitamos definir la matriz identidad.

\begin{apuntebloque}{Definición}
Se define la \textbf{matriz identidad} como
\begin{equation*}
\delta^i{}_j =
\begin{cases}
    1, & \text{si } i = j, \\[6pt]
    0, & \text{si } i \neq j.
\end{cases}
\end{equation*}
\end{apuntebloque}

Escribiendo la definición en forma matricial:
\begin{equation*}
    \delta^i{}_j =
    \begin{pmatrix}
        \delta^1{}_1 & \delta^1{}_2 & \delta^1{}_3 \\
        \delta^2{}_1 & \delta^2{}_2 & \delta^2{}_3 \\
        \delta^3{}_1 & \delta^3{}_2 & \delta^3{}_3
    \end{pmatrix}
    =
    \begin{pmatrix}
        1 & 0 & 0 \\
        0 & 1 & 0 \\
        0 & 0 & 1
    \end{pmatrix}.
\end{equation*}

Donde:
\begin{itemize}
    \item El índice $i$ rotula las \textbf{filas}.
    \item El índice $j$ rotula las \textbf{columnas}.
\end{itemize}

La definición de la matriz identidad es la misma que la definición de la 
métrica Euclidiana. Sin embargo, la definición de la matriz identidad es 
independiente del sistema de coordenadas.  

Usando la definición de la matriz identidad, escribimos el producto de una 
matriz y su inversa en notación tensorial:
\begin{equation*}
    \sum_{k=1}^3 A^i{}_k (A^{-1})^k{}_j
    = \sum_{k=1}^3 (A^{-1})^i{}_k A^k{}_j
    = \delta^i{}_j \,.
\end{equation*}

No se debe olvidar que, para que exista la inversa de una matriz $A$, su 
determinante debe ser distinto de cero:
\begin{equation*}
    \det(A^i{}_j) \neq 0 \,.
\end{equation*}

\begin{apuntebloque}{Definición}
El \textbf{determinante de la métrica} $g_{ij}$ se denota como
\begin{equation*}
    g = \det(g_{ij}) \,.
\end{equation*}
\end{apuntebloque}

Como $g_{ij}$ tiene los índices abajo, entonces la \textbf{matriz inversa} 
tendrá los índices arriba, puesto que la \textbf{matriz identidad} posee 
un índice abajo y uno arriba.

Si $g \neq 0 \;\;\Rightarrow\;\; \exists\, g^{ij}$ tal que
\begin{equation*}
    \sum_{k=1}^3 g^{ik} g_{kj} = \delta^i{}_j,
    \qquad
    \sum_{k=1}^3 g_{jk} g^{ki} = \delta^i{}_j \,.
\end{equation*}


\begin{apuntebloque}{Observación}
La clave de la demostración está en reconocer que:
\begin{equation*}
    (g^{-1})^t \cdot g = I
\end{equation*}

implica que $(g^{-1})^t$ es la inversa de $g$ por la izquierda, y por la unicidad de la inversa, debe ser igual a $g^{-1}$.
\end{apuntebloque}


\begin{apuntebloque}{Proposición}
Si $g_{ij}$ es simétrico (es decir, $g^t = g$), entonces $g^{ij}$ también es simétrico.\\

\textbf{Demostración:}

Sea $g$ una matriz simétrica invertible. Por definición de matriz inversa:
\begin{equation*}
    g \cdot g^{-1} = I \,.
\end{equation*}

Tomando transpuesta en ambos lados:
\begin{equation*}
(g \cdot g^{-1})^t = I^t \,.
\end{equation*}

Por propiedades de la transpuesta:
\begin{equation*}
(g^{-1})^t \cdot g^t = I \,.
\end{equation*}
Dado que $g$ es simétrica ($g^t = g$):
\begin{equation*}
(g^{-1})^t \cdot g = I \,.
\end{equation*}

Esto muestra que $(g^{-1})^t$ es una inversa por la izquierda de $g$. Pero como la inversa de una matriz es única, se concluye que:
\begin{equation*}
(g^{-1})^t = g^{-1} \,.
\end{equation*}

Por lo tanto, $g^{-1}$ es simétrica, lo que en componentes significa:
\begin{equation*}
g^{ij} = g^{ji} \,.
\end{equation*}
\hfill $\square$
\end{apuntebloque}


\begin{apuntebloque}{Ejemplo: Determinante e inversa de la métrica}

\textbf{a) Coordenadas Cartesianas}  \\

La métrica en coordenadas cartesianas es $g_{ij} = \delta_{ij}$, 
la cual tiene determinante uno. Su inversa es
\begin{equation*}
    [\delta^{ij}] =
    \begin{pmatrix}
        1 & 0 & 0 \\
        0 & 1 & 0 \\
        0 & 0 & 1
    \end{pmatrix},
    \qquad
    \delta^{ij} =
    \begin{cases}
        1, & i = j, \\[6pt]
        0, & i \neq j.
    \end{cases}
\end{equation*}

Donde:
\begin{itemize}
    \item El índice $i$ rotula las \textbf{filas}.
    \item El índice $j$ rotula las \textbf{columnas}.
\end{itemize}

La métrica inversa $\delta^{ij}$ es idéntica a la matriz identidad, 
pero no es la matriz identidad: es la \textbf{métrica inversa euclidiana}.

\bigskip

\textbf{b) Coordenadas Esféricas} \\ 

El determinante de la métrica en coordenadas esféricas es
\begin{align*}
    \det(g_{ij}) &=
    \det\!
    \begin{pmatrix}
        1 & 0 & 0 \\
        0 & r^2\sin^2\theta & 0 \\
        0 & 0 & r^2
    \end{pmatrix} \\[4pt]
    &= r^4 \sin^2\theta \,.
\end{align*}

Este determinante es distinto de cero, salvo en $r=0$ y 
$\theta \notin \{ 0,\pi \}$, que son los puntos donde los vectores 
$\vec{e}_\phi$ y $\vec{e}_\theta$ se anulan.  

La métrica inversa en coordenadas esféricas es
\begin{equation*}
    [g^{ij}] =
    \begin{pmatrix}
        1 & 0 & 0 \\
        0 & \dfrac{1}{r^2 \sin^2\theta} & 0 \\
        0 & 0 & \dfrac{1}{r^2}
    \end{pmatrix},
\end{equation*}
la cual diverge en $r=0$ y en $\theta \notin \{0,\pi\}
$.
\end{apuntebloque}

La métrica no solo permite obtener el elemento de línea, 
sino que además contiene información acerca del sistema 
coordenado, como por ejemplo dónde se anulan y/o divergen 
los vectores base coordenados.  

\begin{apuntebloque}{Proposición}
La derivada de $g$ con respecto a la componente $g_{ij}$ satisface
\begin{equation*}
    \frac{\partial g}{\partial g_{ij}} = g \, g^{ij}.
\end{equation*}
\end{apuntebloque}

\newpage

\end{document}
